\input includes/\format

\pdfinfo{
   /Title (Shikatsu Myoki)
   /Creator (TeX)
   /Author (Vit Brunner)
   /CreationDate (20261111233500)
}

~
\title
\parindent=0em
\baselineskip=1em

\vskip 5cm
\centerline{\japanesetitle 死活妙機}
\bigskip
\centerline{Shikatsu Myoki}
\vfill\break

\input includes/\format-motto

\bigskip
{\header Preface}
\medskip
Shikatsu Myoki is a classical life-and-death collection by Honinbo Shusai, first published in 1910. Its practical, game-like problems are aimed at strong dan players.

This booklet contains the 120 main problems, taken from Sensei's Library, omitting the variations, and retaining its four thematic sections. All positions are shown as {\bf black to play}.
\bigskip

\bigskip

\def\chapterentry#1#2#3{%
  \pdfsectionbookmark{#1}{#2}%
  \hbox to \hsize{\strut
    \pdfstartlink attr{/Border [0 0 0]} goto name {problem-#1-1}%
    #1. #2 (#3 problems)\dotfill \chapterpage{#1}%
    \pdfendlink}}
\vbox{
\noindent{\header Sections}
\medskip
\chapterentry{1}{Killing}{48}
\chapterentry{2}{Living}{30}
\chapterentry{3}{Ko}{32}
\chapterentry{4}{Capturing races}{10}
}

\bigskip
\rightline{\it\me}
\rightline{\it November 2026}
\vfill\break

\input includes/\format-gotc
\baselineskip=1.1em
\input generated/\format-shikatsu-myoki

\bye
